\subsection{MORPHISMS}

In this section and the next we shall assume for the sake of simplicity that the species of structures under consideration has only one base set (which is therefore a principal base set). The reader will have no difficulty in extending the definitions and results to the general case.

Let \(\Sigma\) be a species of structures in a theory \(\mathcal{C}\) which is stronger than the theory of sets, and let \(x, y, s, t\) be four distinct letters which are different from the constants of \(\mathcal{C}\) . We recall that the notation \(\mathcal{F}(x, y)\) denotes the set of mappings of \(x\) into \(y\) (Chapter II, § 5, no. 2). Suppose that we are given a term \(\sigma\{x, y, s, t\}\) in \(\mathcal{C}\) which satisfies the following conditions:

(MO \(_{I}\) ) The relation `` \(s\) is a structure of species \(\Sigma\) on \(x\) , and \(t\) is a structure of species \(\Sigma\) on \(y\) '' implies, in \(\mathfrak{C}\), the relation \(\sigma\{x,y,s,t\} \subset \mathcal{F}(x,y)\) .

\((\mathrm{MO}_{\mathrm{II}})\) If, in a theory \(\mathfrak{C}'\) stronger than \(\mathfrak{C}\) , we have three sets \(\mathbf{E}\) , \(\mathbf{E}'\) , \(\mathbf{E}''\) endowed respectively with structures \(\mathscr{S}\) , \(\mathscr{S}'\) , \(\mathscr{S}''\) of species \(\Sigma\) , then the relations \(f \in \sigma \{ \mathbf{E}, \mathbf{E}', \mathscr{S}, \mathscr{S}' \}\) and \(g \in \sigma \{ \mathbf{E}', \mathbf{E}'', \mathscr{S}', \mathscr{S}'' \}\) imply the relation

\[
g \circ f \in \sigma \left\{\mathrm{E}, \mathrm{E} ^ {\prime \prime}, \mathscr {S}, \mathscr {S} ^ {\prime \prime} \right\}.
\]

\((\mathrm{MO}_{\mathrm{III}})\) If, in a theory \(\mathfrak{C}'\) stronger than \(\mathfrak{C}\) , we have two sets \(\mathbf{E}\) , \(\mathbf{E}'\) endowed respectively with structures \(\mathcal{S}\) , \(\mathcal{S}'\) of species \(\Sigma\) , then a bijection \(f\) of \(\mathbf{E}\) onto \(\mathbf{E}'\) is an isomorphism if and only if \(f \in \sigma \left\{\mathbf{E}, \mathbf{E}', \mathcal{S}, \mathcal{S}'\right\}\) and \(f^{-1} \in \sigma \left\{\mathbf{E}', \mathbf{E}, \mathcal{S}', \mathcal{S}\right\}\) .

If \(\Sigma\) and \(\sigma\) are given, the relation \(f \in \sigma\{x, y, s, t\}\) is expressed by saying that \(f\) is a morphism (or a \(\sigma\) -morphism) of \(x\) , endowed with \(s\) , into \(y\) , endowed with \(t\) . If (in a theory \(\mathfrak{G}'\) stronger than \(\mathfrak{G}\) ) \(\mathbf{E}\) and \(\mathbf{E}'\) are two sets endowed with structures \(\mathcal{S}, \mathcal{S}'\) of species \(\Sigma\) , then the term \(\sigma\{\mathbf{E}, \mathbf{E}', \mathcal{S}, \mathcal{S}'\}\) is the set of \(\sigma\) -morphisms of \(\mathbf{E}\) into \(\mathbf{E}'\) .

Examples

\begin{enumerate}
\def\labelenumi{(\arabic{enumi})}
\item
  Take \(\Sigma\) to be the species of order structures and let \(\sigma\{x, y, s, t\}\) denote the set of all mappings \(f\) of \(x\) into \(y\) such that the relation \((u, v) \in s\) implies \((f(u), f(v)) \in t\) . With the notation of Chapter III, § 1, this means that \(u \leqslant v\) implies \(f(u) \leqslant f(v)\) , i.e., that \(f\) is increasing. The verification of axioms \((\mathrm{MO}_{\mathbf{I}}), (\mathrm{MO}_{\mathbf{II}})\) , and \((\mathrm{MO}_{\mathbf{III}})\) is obvious.
\item
  Take \(\Sigma\) to be a species of algebraic structures which has a single (internal) law of composition, everywhere defined (§1, no. 4, Example 2). Let A, A' be two sets endowed with structures of species \(\Sigma\) , and let \(p, p'\) be the composition laws of these two structures. Consider the mappings \(f\) of A into A' such that \(p'(f(x), f(y)) = f(p(x, y))\) for all \(x \in A\) and all \(y \in A\) . These mappings satisfy (MO \(_{I}\) ), (MO \(_{II}\) ), and (MO \(_{III}\) ), and are called homomorphisms of A into A'.
\item
  Take \(\Sigma\) to be the species of topological structures (§1, no. 4, Example 3). Let A, A' be two sets endowed with topologies V, V', respectively. Consider the mappings \(f\) of A into A' such that the
\end{enumerate}

relation \(\mathbf{X}' \in \mathbf{V}'\) implies \(\overline{f}^{1}(\mathbf{X}') \in \mathbf{V}\) (in other words, such that the inverse image of every open set in the topology \(\mathbf{V}'\) is an open set in the topology \(\mathbf{V}\) ). These mappings, which satisfy \((\mathrm{MO}_{\mathbf{I}})\) , \((\mathrm{MO}_{\mathbf{II}})\) , and \((\mathrm{MO}_{\mathbf{III}})\) , are the continuous mappings of \(\mathbf{A}\) into \(\mathbf{A}'\) (with respect to the topologies \(\mathbf{V}\) and \(\mathbf{V}'\) ) (cf.~General Topology, Chapter I, § 2).

Remark. For given species of structures \(\Sigma\) we may have occasion to define various terms \(\sigma \{x, y, s, t\}\) which satisfy the conditions \((\mathrm{MO}_{\mathbf{I}})\) , \((\mathrm{MO}_{\mathbf{II}})\) , and \((\mathrm{MO}_{\mathbf{III}})\) . For example, if \(\Sigma\) is the species of topological structures, with the notation of Example 3 above, a mapping \(f\) of A into \(\mathbf{A}'\) is said to be open if the relation \(\mathbf{X} \in \mathbf{V}\) implies \(f(\mathbf{X}) \in \mathbf{V}'\) (in other words, if the image under \(f\) of every open set is an open set). It is easily checked that the open mappings also satisfy conditions \((\mathrm{MO}_{\mathbf{I}})\) , \((\mathrm{MO}_{\mathbf{II}})\) , and \((\mathrm{MO}_{\mathbf{III}})\) for the species \(\Sigma\) . * Moreover, it can be shown that a continuous mapping is not necessarily open, and that an open mapping is not necessarily continuous. * A given species of structures therefore does not imply a well-defined notion of morphisms.

\pandocbounded{\includegraphics[keepaspectratio]{images/part002_a17f5ac5846160dabe35edfd7ae1b1f71cd54f2b9ab149a24f189ab7736ff1b5.jpg}}

Where order structures, algebraic structures, and topological structures are concerned, it is always to be understood that the morphisms are those which have been defined in the Examples above, unless the contrary is expressly stated.

The condition \((MO_{III})\) and the characterization of bijections (Chapter II, § 3, no. 8, Corollary to Proposition 8) imply the following criterion :

CST8. Let E, E' be two sets, each endowed with a structure of species Σ. Let f be a σ-morphism of E into E' and let g be a σ-morphism of E' into E. If g∘f is the identity mapping of E onto itself, and if f∘g is the identity mapping of E' onto itself, then f is an isomorphism of E onto E', and g is the inverse isomorphism.

\pandocbounded{\includegraphics[keepaspectratio]{images/part002_4a2c9e6ca2e0232c08246aaeec85214d3ec9d06a429c17d5fe3c6f33ea15bbf4.jpg}}

It should be noted that a bijection of E onto E' may be a σ-morphism without the inverse bijection necessarily being a σ-morphism. * For example, a bijective mapping of a topological space A onto a topological space A' can be continuous without the inverse bijection being continuous (General Topology, Chapter I, § 2, no. 1, Remark 1). *

Remark. When a species of structures \(\Sigma\) has several principal base sets \(x_{1},\ldots ,x_{n}\) , and auxiliary base sets \(\mathrm{A}_1,\dots ,\mathrm{A}_m\) , a \(\sigma\) -morphism is a system \((f_1,\dots ,f_n)\) , where \(f_{i}\) is a mapping of \(x_{i}\) into \(y_{i}(1\leqslant i\leqslant n)\) , and these systems of mappings satisfy conditions analogous to (MO \(_{\mathrm{II}}\) ) and (MO \(_{\mathrm{III}}\) ), which the reader may easily state for himself.

